\section{Results}\label{sec:results}

\subsection*{A pre-registered pipeline evaluated in three external directions}

The development pool comprised \pn{n.dev} images (\pn{n.idrid} from IDRiD and
\pn{n.messidor2} from Messidor-2) under a frozen 5-fold split; the external sets were APTOS
2019 (\pn{n.aptos} images) and DDR (\pn{n.ddr} images), and a separate reverse-direction model
was trained on APTOS and tested on Messidor-2 and IDRiD (Fig.~\ref{fig:design},
Table~\ref{tab:cohorts}). The selected model, an EfficientNet-B3 with two ordinal heads,
was chosen on the development pool before any external number was examined. It reached a
development QWK of \pn{sel.dev.qwk} at matched calibration and an external QWK of
\pn{sel.aptos.qwk} (95\% CI \pn{sel.aptos.lo}--\pn{sel.aptos.hi}) on APTOS and \pn{ddr.sel.R1.qwk}
on DDR at its shipped decision rule. On identical images, folds and decoding, the full
pipeline was worth \pn{pipeline.worth} QWK over frozen ImageNet features with a linear ordinal
head (QWK \pn{probe.imagenet}). These numbers are competitive; the remainder of the paper asks
what they do and do not measure.

\subsection*{Ranking transferred to new populations; the mapping from score to grade did not}

Across the \pn{aptos.nmodels} pre-registered development models, the referral probability
$P(\mathrm{grade}\geq2)$ ranked external images almost as well on both corpora: the area under
the ROC curve was \pn{aptos.auroc.min}--\pn{aptos.auroc.max} on APTOS and
\pn{ddr.auroc.min}--\pn{ddr.auroc.max} on DDR. Where those scores fell relative to the shipped
cut-points did not transfer, and it moved in opposite directions on the two populations. Every
model over-graded APTOS (mean predicted-minus-true grade \pn{aptos.R1.mean.min} to
\pn{aptos.R1.mean.max}) and under-graded DDR (\pn{ddr.R1.mean.min} to \pn{ddr.R1.mean.max}). At the
shipped rule (R1) this placed every model at a referable-DR sensitivity of
\pn{aptos.R1.sens.min}--\pn{aptos.R1.sens.max}\% on APTOS but \pn{ddr.R1.sens.min}--\pn{ddr.R1.sens.max}\%
on DDR (specificity \pn{ddr.R1.spec.min}--\pn{ddr.R1.spec.max}\%), and cut-points refitted on the
external labels (R4, two-fold cross-fitted, a diagnostic) raised DDR sensitivity to
\pn{ddr.R4.sens.min}--\pn{ddr.R4.sens.max}\% without retraining.

A referral threshold fitted on development predictions therefore did not hold its operating
point (Fig.~\ref{fig:transfer}). Fitted for the 85\% sensitivity of the superiority endpoint of
the first autonomous pivotal trial\cite{abramoff2018pivotal}, the reference model referred
\pn{aptos.tr.e08.B.sens}\% of referable patients on APTOS but \pn{ddr.tr.e08.B.sens}\% on DDR
(pre-registered O8 \pn{ddr.o8}); fitted for 80\%, it reached \pn{aptos.tr.e08.A.sens}\% and
\pn{ddr.tr.e08.A.sens}\% (gaps \pn{aptos.tr.e08.A.gap} and \pn{ddr.tr.e08.A.gap} points). The
same threshold sits at different points of the operating curve on different populations; a
development-fitted operating point is not a property of the model.

The reverse direction showed the same dissociation. A model trained on APTOS with the same
recipe (out-of-fold QWK \pn{rev.dev.qwk}) under-graded Messidor-2 (mean predicted-minus-true
grade \pn{rev.m2.R1.mean}; pre-registered O6 \pn{rev.o6}) and, at its shipped rule, referred only
\pn{rev.m2.R1.sens}\% of referable eyes (specificity \pn{rev.m2.R1.spec}\%), although it ranked
them with an area under the curve of \pn{rev.m2.auroc}; refitted cut-points raised sensitivity
to \pn{rev.m2.R4.sens}\% (specificity \pn{rev.m2.R4.spec}\%). The direction of miscalibration was
set by the target population, not by the direction of transfer: the development models
under-graded DDR just as the APTOS model under-graded Messidor-2.

\subsection*{The mild-grade collapse on APTOS is a decision-rule failure, but it is corpus-specific}

On APTOS every development model lost the mild grade. At the shipped rule the models recalled
\pn{aptos.R1.mild.min}--\pn{aptos.R1.mild.max}\% of the 370 mild-DR images (Fig.~\ref{fig:rules}a),
against \pn{aptos.devmild.min}--\pn{aptos.devmild.max}\% of mild images in the development pool
at the same rule. Refitting only the four cut-points on half
of the external labels and applying them to the other half (R4m) restored mild recall to
\pn{aptos.R4m.mild.min}--\pn{aptos.R4m.mild.max}\% in every model, and the same held in all five
seeds of the reference configuration (R1 \pn{seeds.aptos.mildR1.lo}--\pn{seeds.aptos.mildR1.hi}\%,
R4m \pn{seeds.aptos.mildR4m.lo}--\pn{seeds.aptos.mildR4m.hi}\%; pre-registered O3 \pn{seeds.o3}).
Over-grading sent mild images across the moderate cut-point; the representation had separated
them.

The collapse did not generalise beyond APTOS. On DDR, where the models under-graded, mild
recall at the shipped rule was \pn{ddr.R1.mild.min}--\pn{ddr.R1.mild.max}\% and external
cut-points raised it only to \pn{ddr.R4m.mild.min}--\pn{ddr.R4m.mild.max}\% (Fig.~\ref{fig:rules}c).
In the reverse direction, where a recent benchmark reported a mild-grade F1 of 0.000 for frozen
encoders\cite{poyrazer2026calibration}, our fine-tuned model recalled \pn{rev.m2.R1.mild}\% of the
\pn{rev.m2.nmild} mild Messidor-2 images at the shipped rule, errors falling to no DR
(\pn{rev.m2.mild.to0}) rather than to moderate (\pn{rev.m2.mild.to2}); the pre-registered
predictions that it would collapse (O4) and that refitted cut-points would at least triple it
(\pn{rev.m2.R4m.mild}\%; O5) were both falsified. A collapsed grade is thus one possible
consequence of a misplaced decision rule, not a general property of external DR grading.

\subsection*{Deployable rules did not repair the failures, and models differed more than rules}

Rules fitted on development data did not transfer either. Cut-points chosen to maximise
development QWK (R2) lowered external QWK for \pn{aptos.r2worse} of the \pn{aptos.nmodels}
models on APTOS and left mild recall at \pn{aptos.R2.mild.min}--\pn{aptos.R2.mild.max}\%; a referral
cut fixed to give 90\% development sensitivity (R3) behaved likewise. We had pre-registered the
prediction (O7) that the choice among these deployable rules would move mild recall and referable
sensitivity more than the choice among models. It was falsified for both metrics on both corpora:
on APTOS the median within-model range across R1--R3 was \pn{aptos.o7.mildrecall.within} points
of mild recall against a between-model range of \pn{aptos.o7.mildrecall.between} at R1, and
\pn{aptos.o7.refsens.within} against \pn{aptos.o7.refsens.between} points of sensitivity; on DDR
\pn{ddr.o7.mildrecall.within} against \pn{ddr.o7.mildrecall.between} and \pn{ddr.o7.refsens.within}
against \pn{ddr.o7.refsens.between}. Models trained on the same data thus differed materially at
their shipped rules, and no rule fitted without external labels removed the population shift.
The ranking of models was nonetheless rule-dependent: the Spearman correlation between rankings
by external QWK at R1 and at the external-label rule R4 was \pn{aptos.rho.R1vsR4} on APTOS and
\pn{ddr.rho.R1vsR4} on DDR, and the development ranking predicted the APTOS ranking with
$\rho=$ \pn{aptos.rho.devmatchedvsR1}.

\subsection*{The grading metric cannot see the referral operating point}

QWK, the field's summary statistic, was nearly invariant to the decision that determines who
is referred. Holding the other cut-points at their QWK-optimal values and sweeping only the
referral cut (Fig.~\ref{fig:frontier}), every operating point within 0.01 of the maximum QWK
for the selected model on APTOS (\pn{aptos.band.e11full.qwkmax}) spanned referable sensitivity
from \pn{aptos.band.e11full.sens.lo}\% to \pn{aptos.band.e11full.sens.hi}\% and specificity from
\pn{aptos.band.e11full.spec.lo}\% to \pn{aptos.band.e11full.spec.hi}\%: between
\pn{aptos.band.e11full.miss.lo} and \pn{aptos.band.e11full.miss.hi} of
\pn{aptos.band.e11full.nref} referable patients missed, with QWK unable to tell them apart. On
DDR the corresponding band spanned \pn{ddr.band.e11full.sens.lo}--\pn{ddr.band.e11full.sens.hi}\%
sensitivity, between \pn{ddr.band.e11full.miss.lo} and \pn{ddr.band.e11full.miss.hi} of
\pn{ddr.band.e11full.nref} referable patients missed. Optimising for the grade the metric under-weights carried its own cost: cut-points
chosen to maximise macro-averaged recall recovered the mild grade but cut the reference model's
APTOS referable sensitivity from \pn{aptos.e08.R1.sens}\% to \pn{aptos.e08.R4m.sens}\%, because
mild DR is not a referable grade. Macro-recall is therefore reported here as a diagnostic and
never optimised.

\subsection*{The same pattern in a published model family}

Because ORDER-DR\cite{sheng2026orderdr} released image-level class probabilities for its models
on both its APTOS validation split and the \pn{odr.n} Messidor-2 images, we could apply the
same rules to an independently trained model family in the published direction, without any
weights. Using the authors' own fusion and cut-points we reproduced their reported external
result exactly (QWK \pn{odr.odr.R2.qwk} $\pm$ \pn{odr.odr.R2.qwksd} over three seeds;
\pn{odr.odr.R2.pred1} images predicted mild against \pn{odr.nmild} true, precision
\pn{odr.odr.R2.prec}; Table~\ref{tab:orderdr}). At argmax, the four released models recalled
\pn{odr.R1.mild.min}--\pn{odr.R1.mild.max}\% of mild images and \pn{odr.R1.sens.min}--\pn{odr.R1.sens.max}\%
of referable ones. With cut-points fitted on external labels (R4), all four converged to QWK
\pn{odr.R4.qwk.min}--\pn{odr.R4.qwk.max}: the 0.07 QWK advantage reported for the fused model
over its best branch was a decision-rule effect. Refitting for macro-recall roughly doubled
referable sensitivity (\pn{odr.R4m.sens.min}--\pn{odr.R4m.sens.max}\%) and raised mild recall to
\pn{odr.R4m.mild.min}--\pn{odr.R4m.mild.max}\%. Mild F1, however, stayed at or below
\pn{odr.R4m.f1.max}: in this direction part of the mild failure is representational, and the
decomposition, not either extreme, is the result.

\subsection*{Matched calibration changes model comparisons, and seeds set the noise floor}

Comparisons between models were as sensitive to the decision rule as single-model metrics
(Table~\ref{tab:record}). On APTOS the selected EfficientNet-B3 led the DenseNet121 reference
by \pn{rec.bb.R1.qwk} QWK \pn{rec.bb.R1.qwkci} at the shipped rule; with both models'
cut-points refitted the same way the difference became \pn{rec.bb.R4.qwk} \pn{rec.bb.R4.qwkci},
and per-grade differences changed sign across operating points (moderate recall
\pn{rec.bb.R1.mod}, \pn{rec.bb.R4.mod} and \pn{rec.bb.R4m.mod} points under R1, R4 and R4m;
severe \pn{rec.bb.R1.sev}, \pn{rec.bb.R4.sev}, \pn{rec.bb.R4m.sev}). On the development pool,
linear-probe-then-fine-tune appeared to cost \pn{rec.lpft.ship} accuracy points
\pn{rec.lpft.shipci} at shipped cut-points and gained \pn{rec.lpft.match} \pn{rec.lpft.matchci}
at matched ones. The rule is not a one-way filter: the development advantage of EfficientNet-B3
over its DenseNet121 control (\pn{iv.effnet.ship} shipped) survived matched calibration
(\pn{iv.effnet.diff} \pn{iv.effnet.ci}), and proliferative-DR recall favoured EfficientNet-B3
at every operating point.

Five seeds of the reference configuration set the scale of these effects. Development QWK at
matched calibration varied with a standard deviation of \pn{seeds.dev.sd} across seeds
(range \pn{seeds.dev.lo}--\pn{seeds.dev.hi}), and APTOS QWK at the shipped rule with
\pn{seeds.aptos.sd}. Between models that differed only by seed, \pn{seeds.aptos.sigR1} of
\pn{seeds.aptos.ntests} per-grade recall comparisons on APTOS were ``significant'' at the
shipped rule (largest difference \pn{seeds.aptos.maxR1} points; pre-registered O2
\pn{seeds.o2}), and \pn{seeds.aptos.sigR4} at matched external cut-points; on DDR the counts
were \pn{seeds.ddr.sigR1} and \pn{seeds.ddr.sigR4} of \pn{seeds.ddr.ntests}. An image-level
bootstrap captures sampling of test images, not the variance of training, so per-grade
differences of the size the literature attributes to architecture arise between identical
architectures. A difference between configurations is treated here as beyond seed noise only
if it exceeds $2\sqrt{2}$ seed standard deviations of matched development QWK
(\pn{seeds.yardstick}; Fig.~\ref{fig:forest}).

\subsection*{About 200 local labels are enough to recalibrate safely}

If cut-points must be refitted locally, the practical question is how many labelled local
images that requires. Fitting cut-points on $n$ randomly drawn APTOS images and evaluating on the
rest (Fig.~\ref{fig:recal}), \pn{recal.25.harm}\% of draws at $n=25$ and \pn{recal.100.harm}\%
at $n=100$ ended worse than the shipped rule, although the mean gain was positive at every size.
At $n=200$ recalibration recovered \pn{recal.200.gain} macro-recall points \pn{recal.200.ci}
with a \pn{recal.200.harm}\% probability of harm, and at $n=400$ the probability of harm was
\pn{recal.400.harm}\%. A recommendation judged by its mean alone would have endorsed a
procedure that harms roughly one adopter in four at small $n$.

\subsection*{Interventions on the representation, and a supervision ceiling for macular oedema}

Of \pn{iv.n} pre-registered interventions on the model itself, \pn{iv.nfalsified} failed their
pre-specified criterion at matched calibration (Fig.~\ref{fig:forest}; Supplementary Table~S3).
Only the EfficientNet-B3 backbone improved development QWK; a six-model ensemble gained on the
development pool and \pn{ensemble.diff} \pn{ensemble.ci} on APTOS. The retinal foundation model
showed the clearest dissociation between representation and outcome: frozen, its features
carried more DR signal than ImageNet features under an identical linear probe
(\pn{probe.retfound} versus \pn{probe.imagenet}, \pn{probe.diff} \pn{probe.ci}); fine-tuned with
linear-probe-then-fine-tune at its native 224 pixels, it fell below a DenseNet121 pretrained on
EyePACS at the same resolution by \pn{iv.retfound.diff} \pn{iv.retfound.ci} over
\pn{iv.retfound.folds} folds (pre-registered O9 \pn{retfound.o9}). Externally (a secondary
analysis, outside the pre-registered model set) its two failures had different causes. On APTOS
it ranked referable eyes as well as the convolutional models (area under the curve
\pn{rf.aptos.auroc}) and most of its QWK deficit was the decision rule: \pn{rf.aptos.R1.qwk} at
its shipped rule against \pn{rf.aptos.R4.qwk} with refitted cut-points. On DDR its ranking itself
degraded (area under the curve \pn{rf.ddr.auroc}, QWK \pn{rf.ddr.R4.qwk} even with refitted
cut-points), a representational failure that recalibration cannot repair.

No intervention moved the three-grade DME head. A crop centred on the fovea by a localiser
validated before use (median error \pn{fovea.median} disc diameters), which encodes the
clinical definition of the grade directly, changed DME QWK by \pn{iv.macula.diff}
\pn{iv.macula.ci}. The constraint is supervision: only \pn{dme.n1} of the \pn{n.idrid}
three-grade images (\pn{dme.mid.pct}\%) carry the middle grade, a proportion that a
purpose-built multimodal dataset reproduces (\pn{lit.p6.mid}\%)\cite{tang2026mmrdr}, and a
deployed commercial system reached a DME sensitivity of \pn{lit.p12.dme.sens}\% ($\kappa$
\pn{lit.p12.dme.kappa}) on the same images on which its referable-DR $\kappa$ was
\pn{lit.p12.dr.kappa}\cite{duggal2025realworld}. Nor can the grade be derived from published
annotations: applying the definition (nearest hard exudate within one disc diameter of the
fovea) to IDRiD's pixel-level exudate masks matched the expert grade in \pn{parta.exact}\% of
\pn{parta.n} images, against \pn{parta.constant}\% for a constant predictor, and failed on two
of the \pn{parta.informative} images where the test was informative, each time on a single
small exudate lying just inside the threshold (Supplementary Note~2).
